\subsection{Spectral projectors}

Let \(u\) be a normal partial operator on a complex Hilbert space E. Let A be a universally measurable subset of \(\mathrm{Sp}(u)\) and \(\varphi_{\mathrm{A}}\) its characteristic function. Since \(\varphi_{\mathrm{A}}\) is bounded and satisfies \(\varphi_{\mathrm{A}}^{2} = \varphi_{\mathrm{A}}\) , the endomorphism \(\varphi_{\mathrm{A}}(u)\) of E is an orthogonal projector of E. It is called the spectral projector of \(u\) defined by A. We denote \(p_{\mathrm{A}} = \varphi_{\mathrm{A}}(u)\) . We have \(p_{\emptyset} = 0\) and \(p_{\mathrm{Sp}(u)} = 1_{\mathrm{E}}\) .

Let A be a universally measurable subset of C. The spectral projector of u defined by A is the spectral projector \(p_{\mathrm{Sp}(u)\cap\mathrm{A}}\) . It is also denoted simply \(p_{A}\) . For every function \(f \in \mathcal{L}_{\mathrm{u}}(\mathrm{Sp}(u))\) , for all \(x \in E\) and every \(y \in \mathrm{dom}(f(u))\) , we have the formula

\[
\langle p _ {\mathrm{A}} (x) \mid f (u) y \rangle = \int_ {\operatorname{Sp} (u) \cap \mathrm{A}} f d \mu ,\tag{10}
\]

where \(\mu\) is the spectral measure of \((x,y)\) relative to \(u\) (formula (7), p. 271). Let A and B be universally measurable subsets of \(\mathrm{Sp}(u)\) . Since \(\varphi_{\mathrm{A}}\varphi_{\mathrm{B}} = \varphi_{\mathrm{A}\cap \mathrm{B}}\) , we have \(p_{\mathrm{A}} \circ p_{\mathrm{B}} = p_{\mathrm{A}\cap \mathrm{B}}\) (prop. 5 of IV, p. 275, \(c\) )). In particular, if A and B are disjoint, the images of \(p_{\mathrm{A}}\) and of \(p_{\mathrm{B}}\) are orthogonal.

PROPOSITION 8. --- Let \((\mathrm{A}_{i})_{i\in\mathrm{I}}\) be a countable family of pairwise disjoint universally measurable subsets of \(\mathrm{Sp}(u)\) and let \(p_{i}\) be the spectral projector of u defined by \(A_{i}\) . The union A of the sets \(A_{i}\) is a universally measurable subset of \(\mathrm{Sp}(u)\) and the series \(\sum_{i}p_{i}\) converges to \(p_{A}\) in \(\mathcal{L}(\mathrm{E})\) endowed with the topology of pointwise convergence.

The set A is universally measurable by INT, IV, p. 177, § 5, n° 4, cor. 2. The series \(\sum_{i} \varphi_{\mathrm{A}_{i}}\) converges simply to \(\varphi_{\mathrm{A}}\) and its partial sums are bounded by 1. The assertion therefore follows from prop. 6 of IV, p. 276.

PROPOSITION 9. --- Let A be a closed subset of Sp(u). Let \(\varphi_{\mathrm{A}}\) be the characteristic function of A and \(p_{\mathrm{A}} = \varphi_{\mathrm{A}}(u)\) the spectral projector of u defined by A. Denote by \(\mathrm{E}_{\mathrm{A}}\) the image of \(p_{\mathrm{A}}\) . This is a closed subspace of E.

\begin{enumerate}
\def\labelenumi{\alph{enumi})}
\item
  The subspace \(E_{A}\) is the set of \(x \in E\) such that the support of the spectral measure of x relative to u is contained in A;
\item
  If A is bounded in \(\mathbf C\), then \(E_{A}\) is contained in the domain of u;
\item
  For all x belonging to the domain of u, we have \(p_{\mathrm{A}}(x) \in \operatorname{dom}(u)\) and \(u(p_{\mathrm{A}}(x)) \in \mathrm{E}_{\mathrm{A}}\) in particular \(u(x) \in \mathrm{E}_{\mathrm{A}}\) if \(x \in \operatorname{dom}(u) \cap \mathrm{E}_{\mathrm{A}}\) .
\end{enumerate}

For \(x\) in E, we denote by \(\mu_x\) the spectral measure of \(x\) relative to \(u\) .

Let us prove \(a\) ). Let \(x\) be an element of \(\mathrm{E}_{\mathrm{A}}\) . Let \(z \in \mathbf{C} - \mathbf{A}\) and let U be a relatively compact open neighborhood of \(z\) which does not meet A. For every function \(f \in \mathcal{K}(\mathbf{C})\) with support contained in U, we have

\[
\int_ {\mathrm{Sp} (u)} f \mu_ {x} = \langle x | f (u) x \rangle = \langle p _ {\mathrm{A}} (x) | f (u) x \rangle = \int_ {\mathrm{Sp} (u) \cap \mathrm{A}} f \mu_ {x} = 0
\]

by formula (10). This means that \(z\) does not belong to the support of \(\mu_x\) . Consequently, the support of \(\mu_x\) is contained in A.

Conversely, let \(x \in E\) such that the support of \(\mu_{x}\) is contained in A. It follows that

\[
\langle x \mid p _ {\mathrm{A}} (x) \rangle = \int_ {\operatorname{Sp} (u) \cap \mathrm{A}} \mu_ {x} = \int_ {\operatorname{Sp} (u)} \mu_ {x} = \langle x \mid x \rangle ,
\]

\((loc. \ cit.) \ \text{therefore} \ \|p_{\mathrm{A}}(x)-x\|^{2} = \|p_{\mathrm{A}}(x)\|^{2} - \|x\|^{2} \leqslant 0 \ \text{and consequently } p_{\mathrm{A}}(x) = x, \ \text{that is, } x \in \mathrm{E}_{\mathrm{A}}.\)

Assertion b) follows from a) and the remark in IV, p. 273.

Let us prove \(c\) ). The domain of \(u\) is the set of \(x \in \mathrm{E}\) such that the identity function of \(\operatorname{Sp}(u)\) belongs to \(\mathcal{L}^2(\operatorname{Sp}(u), \mu_x)\) (loc. cit.). Since \(\mu_{p_{\mathrm{A}}(x)} = \varphi_{\mathrm{A}} \cdot \mu_x\) for \(x \in \mathrm{E}\) (corollary of Prop. 5 of IV, p. 275), we have \(p_{\mathrm{A}}(x) \in \operatorname{dom}(u)\) if \(x \in \operatorname{dom}(u)\) .

Let \(x \in \operatorname{dom}(u)\) and \(y = p_{\operatorname{A}}(x)\) ; we have \(y \in \operatorname{dom}(u) \cap \operatorname{E}_{\operatorname{A}}\) according to \(c\) ). The spectral measure of \(u(y)\) relative to \(u\) is \(|\operatorname{Id}_{\operatorname{Sp}(u)}|^2 \cdot \mu_y\) (loc. cit.). Since \(\mu_y\) is supported in A, the same holds for \(\mu_{u(y)}\) , hence \(u(y) \in \operatorname{E}_{\operatorname{A}}\) according to \(a\) .

COROLLARY. --- Let \(\lambda \in \mathrm{Sp}(u)\) The image of \(p_{\{\lambda\}}\) is the eigensubspace of \(u\) relative to \(\lambda\) .

Let \(x \in \operatorname{dom}(u)\) . Denote by \(\mu\) (resp. \(\nu\) ) the spectral measure of the element \(x\) relative to \(u\) (resp. the spectral measure of \((x, u(x))\) relative to \(u\) ). We have \(\nu = \operatorname{Id}_{\operatorname{Sp}(u)} \cdot \mu\) (corollary of Prop. 5 of IV, p. 275). If \(u(x) = \lambda x\) , we also have \(\nu = \lambda \mu\) , whence the equality \(\operatorname{Id}_{\operatorname{Sp}(u)} \cdot \mu = \lambda \mu\) . This implies that the support of \(\mu\) is contained in \(\{\lambda\}\) (cf. INT, V, p. 46, § 5, no. 3, cor. 2) and therefore \(x\) belongs to the image of \(p_{\{\lambda\}}\) (prop. 9, a)).

Conversely, suppose that \(x\) belongs to the image \(\mathrm{E}_{\lambda}\) of \(p_{\{\lambda\}}\) . Then \(x\) belongs to \(\operatorname{dom}(u)\) and \(u(x)\) also belongs to \(\mathrm{E}_{\lambda}\) (loc. cit., \(b\) )). Since \(\varphi_{\{\lambda\}} \cdot (\operatorname{Id}_{\operatorname{Sp}(u)} - \lambda) = 0\) , we have the relation \(p_{\{\lambda\}} \circ (u - \lambda 1_{\mathrm{E}}) \subset 0\) (prop. 5 of IV, p. 275, \(c\) )), whence \(0 = p_{\{\lambda\}}(u(x)) - \lambda p_{\{\lambda\}}(x) = u(x) - \lambda x\) .

PROPOSITION 10. — Let \(\lambda \in \mathbf{C}\) . We have \(\lambda \in \mathrm{Sp}(u)\) if and only if, for every open neighborhood V of \(\lambda\) in \(\mathbf{C}\) , the spectral projector \(p_{\mathrm{V}}\) of u relative to V is nonzero.

If \(\lambda \notin \mathrm{Sp}(u)\) , then there exists an open neighborhood \(V\) of \(\lambda\) in \(\mathbf{C}\) which does not meet \(\mathrm{Sp}(u)\) , and then \(p_{\mathrm{V}} = p_{\varnothing} = 0\) .

Conversely, suppose that there exists an open neighborhood V of \(\lambda\) in \(\mathbf{C}\) such that \(p_{V} = 0\) . Let c > 0 be such that the disk with center \(\lambda\) and radius c is contained in V.

Let \(x \in \mathrm{dom}(u)\) and let \(\mu_x\) the spectral measure of \(x\) relative to \(u\) . Since \(\mu_x(\mathrm{V}) = \langle x | p_{\mathrm{V}}(x) \rangle = 0\) , one computes

\[
\begin{array}{r l} & {\int_ {\mathbf {C}} | z - \lambda | ^ {2} d \mu_ {x} (z) = \int_ {\mathbf {C - V}} | z - \lambda | ^ {2} d \mu_ {x} (z)} \\ & {\qquad \geqslant c ^ {2} \int_ {\mathbf {C - V}} d \mu_ {x} (z) = c ^ {2} \int_ {\mathbf {C}} d \mu_ {x} (z) = c ^ {2} \| x \| ^ {2}.} \end{array}
\]

But, on the other hand, we have

\[
\| u (x) - \lambda x \| ^ {2} = \int_ {\mathbf {C}} | z - \lambda | ^ {2} d \mu_ {x} (z) = \| u ^ {*} (x) - \overline {{\lambda}} x \| ^ {2}
\]

(formula (8), p. 272), whence \(\|u(x)-\lambda x\|\geqslant c\|x\|\) and \(\|u^{*}(x)-\overline{\lambda}x\|\geqslant c\|x\|\) . It follows that \(\lambda\) belongs to the resolvent set of u (lemma 5 of IV, p. 248). This concludes the proof.

COROLLARY. — Let A be an open set in \(\mathbf{C}\) and \(n \in \mathbf N\) . If A contains n elements of \(\mathrm{Sp}(u)\) , then the dimension of the image of the spectral projector \(p_{A}\) of u relative to A is at least equal to n.

Let \(\lambda_1, \ldots, \lambda_n\) be distinct elements of the spectrum of \(u\) belonging to A. There exists a family \((\mathrm{V}_i)_{1 \leqslant i \leqslant n}\) of pairwise disjoint open sets of \(\mathbf{C}\) such that \(\lambda_i \in \mathrm{V}_i\) for \(1 \leqslant i \leqslant n\) . Let B be the union of the sets \(\mathrm{V}_i\) . The image of \(p_{\mathrm{A}}\) contains the image of the spectral projector \(p_{\mathrm{B}}\) ; moreover, since \(p_{\mathrm{B}}\) is the sum of the projectors \(p_{\mathrm{V}_i}\) , and since the image of \(p_{\mathrm{V}_i}\) is orthogonal to the image of \(p_{\mathrm{V}_j}\) for all \(i \neq j\) , the result follows from Prop. 10.

